\documentclass[../../../include/open-logic-section]{subfiles}

\begin{document}

\olfileid{sth}{story}{predicative}	

\olsection{پريديکاتيف او امپريديکاتيف}

د رسل سټ $R$ د $\Setabs{x}{x \notin
x}$ له لارې تعريف شوے و۔ په بشپړه تشريح، $R$ به هغه سټ وے
چې ټول او يوازې هغه سټونه لري چې خپل ځان د غړي په توګه نۀ
لري۔ د نسخې څرګندونه (OLSTH-001): منجمده سرچينه په دې نثري
تشريح کښې اضافي نفي ليکي، چې د خپل ليکل شوي سټ له فورمول
سره يې معنا سرچپه کوي؛ دلته د عين فورمول معنا بيان شوې ده۔
نو د $R$ په تعريف کښې پر داسې دامن کمي کول کوو چې $R$ به
پکښې شامل وے (که وجود يې لرلے)۔

دا يو \emph{امپريديکاتيف} تعريف دے۔ په عمومي توګه، ويلے شو
چې يو تعريف امپريديکاتيف دے که او يوازې که پر داسې دامن
کمي کول کوي چې تعريفېدونکے شے پکښې شامل وي۔

د تناقضونو له راولاړېدو وروسته، وايټ هېډ، رسل، پوانکارې او
وايل دغسې امپريديکاتيف تعريفونه د «بدې دايرې» په توګه رد کړل:
\begin{quote}
  د هغو تناقضونو شننه چې ترې بايد بچ شو، ښيي چې دا ټول له
  يو ډول بدې دايرې راولاړېږي۔ دغه بدې دايرې له دې فرضه
  پيدا کېږي چې د شيانو يوه ټولګه داسې غړي لرلے شي چې يوازې
  د ټولګې د ټول په وسيله تعريف کېدلے شي[\ldots. \textparagraph]

  هغه اصل چې د نامشروع کليتونو څخه مو بچ کوي، داسې بيانېدلے
  شي: «هر څه چې د يوې ټولګې \emph{ټول} رانغاړي، بايد د هغې
  ټولګې يو غړے نۀ وي»؛ يا په سرچپه توګه: «که يوه ټاکلې ټولګه
  ټول لرلے، نو داسې غړي به يې لرلے چې يوازې د هغه ټول له مخې
  تعريفېدلے شي، نو دغه ټولګه ټول نۀ لري۔» دې ته به «د بدې
  دايرې اصل» وايو؛ ځکه د نامشروع کليتونو په فرض کښې له
  راګډو بدو دايرو څخه مو بچ کوي۔
  \citep[p.~37]{WhiteheadRussell1910}
\end{quote}
که د هغوي په څېر \emph{د بدې دايرې اصل} ومنو، نو ښايي د
ساده ادراک زيانمنه شېما (د \olref[sth][story][rus]{sec})
په داسې څه بدله کړو۔ د نسخې څرګندونه (OLSTH-002): منجمده
سرچينه دلته د دې اصل ردول ليکي؛ مخکينے نقل او راتلونکے
پريديکاتيف بديل د اصل منل غواړي، چې دلته روښانه شوي دي:

\begin{defish}
\emph{پريديکاتيف ادراک۔} د هر فارمول $\phi$ لپاره چې يوازې پر سټونو کمي کول کوي: سټ$^\prime$ $\Setabs{x}{\phi(x)}$ وجود لري۔
\end{defish}

تر څو چې سټونه$^{\prime}$ سټونه نۀ وي، تناقض به نۀ پيدا کېږي۔

له بده مرغه، پريديکاتيف ادراک ډېر \emph{جامع} نۀ دے۔ ځکه
دا له نويو شيانو، سټونو$^\prime$، سره مو اشنا کوي۔ نو بايد
هغه فارمولونه وڅېړو چې پر سټونو$^\prime$ کمي کول کوي۔ که دا
تل يو سټ$^\prime$ راکوي، نو د رسل تناقض به بيا راولاړ شي؛
يوازې د هغو ټولو سټونو$^\prime$ د سټ$^\prime$ په څېړلو چې
خپل ځان د غړي په توګه نۀ لري۔ نو د همدغه فکر په تعقيب بايد
ووايو چې پر سټونو$^\prime$ کمي کول کوونکے فارمول يو اړوند
سټ$^{\prime\prime}$ راکوي۔ بيا به سټونه$^{\prime\prime\prime}$،
سټونه$^{\prime\prime\prime\prime}$، او داسې نور پکار وي۔ د
پرايمونو د ډېرېدو د مخنيوي لپاره به دا اسانه وي چې د
سټونو$_0$، سټونو$_1$، سټونو$_2$، سټونو$_3$، سټونو$_4$،\ldots
په توګه يې وګڼو۔ دا به د ډولونو (ساده) تيورۍ ته د ننوتلو
يوه لاره راکړي۔

پر داسې تيورۍ څو څرګند اعتراضونه شته (که څه هم دا څرګنده
نۀ ده چې \emph{پرېکنده} اعتراضونه دي)۔ په لږو ټکو: له پيدا
شوې تيورۍ سره کار کول ستونزمن دي؛ د شيانو د بېلو ډولونو په
فرض کولو کښې ډېرښت کوي؛ او په پاے کښې دا روښانه نۀ ده چې
امپريديکاتيف تعريفونه ان \emph{دومره بد} دي۔

د وروستي ټکي د روښانولو لپاره د \citeauthor{Ramsey1925}
دا خبره وګورئ:
\begin{quote}
  يو سړي ته په يوه ډله کښې د تر ټولو لوړ سړي په توګه اشاره
  کولے شو؛ په دې توګه يې د هغه کليت په وسيله پېژنو چې هغه
  پخپله يې غړے دے، بې له دې چې کومه بده دايره پکښې وي۔ \citep{Ramsey1925}
\end{quote}
د رمزي خبره دا ده چې «په ډله کښې تر ټولو لوړ سړے» يو
امپريديکاتيف تعريف \emph{دے}؛ خو په څرګنده بېخي روا دے۔

ځواب دا کېدلے شي چې په دې حالت کښې تر ټولو لوړ کس په
\emph{پريديکاتيف} وسيله هم بېلولے شو۔ د بېلګې په توګه، ښايي
هغه سړي ته يوازې ګوته ونيسو۔ نو پر امپريديکاتيف تعريفونو
اعتراض به بايد په څرګنده هغو شيانو پورې محدود وي چې
\emph{يوازې} په امپريديکاتيف ډول بېلېدلے شي۔ خو بيا هم لا
زياته تشريح پکار ده چې ولې دغسې «لازمي امپريديکاتيف‌والے»
دومره بد دے۔\footnote{د نورو معلوماتو لپاره \citet{Linnebo2010} وګورئ۔}

دا منو چې امپريديکاتيف تعريفونه ډېرې بدې پايلې لري، که غواړو
زمونږ تعريفونه د يو شي د \emph{جوړولو} د نسخې په څېر څه
راکړي۔ ځکه له يوه امپريديکاتيف تعريف سره به په رښتيا په بده
دايره کښې ايسار يو: د امپريديکاتيف ډول ټاکل شوي شي د جوړولو
لپاره به \emph{لومړے} ټول شيان جوړوو (هماغه امپريديکاتيف ډول
ټاکل شوے شے هم)، ځکه دغه شے د ټولو شيانو له مخې ټاکل شوے
دے؛ نو د امپريديکاتيف ډول ټاکل شوي شي تر جوړولو مخکې به
هماغه شے جوړوو۔ خو بيا هم دا يوازې هغه وخت د «لازمي
امپريديکاتيف» ډول ټاکل شويو سټونو پر وړاندې يو جدي اعتراض
دے چې سټونه داسې شيان وګڼو چې مونږ يې \emph{جوړوو}۔ او
مونږ يې (ښايي) داسې نۀ ګڼو۔

نو، د ښه يا بد لپاره، هغه لاره چې عامه شوه، د
(ام)پريديکاتيف‌والي په باب سخت دريځ نۀ نيسي۔ بلکې هغه څه
رانغاړي چې اوس د تجمعي-تکراري لارې په توګه ګڼل کېږي۔ په
پاے کښې به دا اجازه راکړي چې خپل سټونه په «پړاوونو» کښې
طبقه‌بندي کړو، \emph{يو څه} لکه پريديکاتيف لاره چې شيان په
سټونو$_0$، سټونو$_1$، سټونو$_2$، \ldots کښې طبقه‌بندي کوي؛
خو د هغو تر منځ به د نوعيت هيڅ توپير فرض نۀ کړو۔

\end{document}
